
\chapter{Further Reading}
What we discuss in the lectures is just a whistle stop tour of many interesting areas of physics. If you are interested in learning more about particular topics or exploring the world of physics beyond what you met in this course I will provide links to further resources and suggestions for further reading.\\

The material linked to in this section may be more advanced than what we discussed during the module so do not feel that you need to understand everything here. The links are grouped together to correspond with the week of the course that they most closely relate to.\\

A nice complement to this module focussing on the concepts behind the physical phenomena rather than the the computational details can be found in the \textbf{Physics playlist} on the Crash Course YouTube channel linked to \href{https://www.youtube.com/playlist?list=PL8dPuuaLjXtN0ge7yDk_UA0ldZJdhwkoV}{here}. The videos in the playlist goes beyond the content of this module but they should all be at an accessible level for you.\\

Another great general resource for physics is the website \href{http://hyperphysics.phy-astr.gsu.edu/hbase/index.html}{Hyper Physics}.  This is a great resource to find  out more about basically any topic in physics, from the mechanics and circuits problems that we have discussed in this module, through to thermodynamics and statistical physics.\\

If you want a more advanced book about \textbf{classical mechanics}, which is what physicists call the sort of physics that we have been studying in this module, then \citep{susskind_classical} can be an enjoyable read. It is based on a lecture course that Leonard Susskind gave at Stanford university aimed at the general public. To get the most out of this book you would need to be comfortable using calculus. However, you can still get a lot out of it even without going through all of the computations.  The book is split into chapters corresponding to the different lectures within the course with interludes between some of the lectures providing some extra background information. Lectures 1, 2, 3, and 5 correspond most closely to what we have discussed in this module, though there is not a perfect overlap. I would recommend taking a look at this book if you are finding this module too easy and want to explore more physics. Personally, I think it can make for some nice bedtime reading.\\

A good book for general problem solving tips is \textbf{Feynman's Tips on Physics}, \citep{feynmantips}, which is a companion book to the famous ``big red books'' of the Feynman lectures, \citep{feynmanv1,feynmanv2,feynmanv3}, and contains lots of tips on how to solve physics problems. Much of the content is at the same level as this module, though the scope of the book goes beyond what we have discussed here.  The full Feynman lectures, available in the library, can also be a nice complement to this module as they cover all of the topics that Richard Feynman thought were essential for a first year physic course to cover. They go much more in depth than we needed to in this module. However, they contain some wonderful explanations and if you are going to study any more physics later in your degree then you should have a look at these.\\

The \textbf{Fun to Imagine} \href{https://www.youtube.com/playlist?app=desktop&list=PL2D30B1DEFFDA0310}{playlist} on Youtube is a collection of videos of Feynman giving explanations of various physical phenomena and expanding on his philosophy for understanding physics.  The videos are all quite short and are clips taken from a BBC programme. I think that they are great fun to listen to, and contain lots of great explanations.

\section{What is Physics Extra Reading}
A standard reference for powers of ten and scales of measurement is the \href{https://www.youtube.com/watch?v=0fKBhvDjuy0}{Powers of Ten} video from 1977 that we saw in the lectures.\\

For more on the question of ``What is Physics?'' you can look at the \href{https://www.tntech.edu/cas/physics/aboutphys/about-physics.php}{Tennessee Tech} page about physics and its different subfields. You can also look at the \href{https://www.iop.org/about}{about} page on the IOP website. The IOP (Institute of Physics) is the learned society and professional body responsible for the promotion of physics in the UK and Ireland.

\section{Kinematics in 1D Extra Reading}
For more details on motion in one dimension have a look at any of the textbooks on the course reading list. The material in \citep{breithaupt2016aqa} and \citep{breithaupt2016aqa_AS} are at exactly the same level as this module and contain explanations and examples of the kinematic equations.\\

If you want more detail at a slightly more advanced level then I recommend chapters 1 and 2 of \citep{YandF2019}. The book \citep{mansfield2020understanding} has some more advanced examples where you need to use ideas from calculus to solve problems from mechanics. If you have already done A-level mathematics and you want to see how it enables you to understand more complicated problems from Physics, then I recommend having at look at this.\\


\section{Kinematics in 2D Extra Reading}
The resources for this week are the same as last week, again if you want to explore mechanics further it is useful to have had some exposure to calculus. Chapter 3 of \citep{YandF2019} is an especially good resource as it starts around the same level as this module before discussing some more complicated problems including motion in three dimensions. \\

There are lots of extra examples and exercise on both one dimensional and two dimensional kinematics in \citep{sadler1996understanding}. Some of the tutorial problems come from this book so it is a great source of problems at the right difficulty level for this module. 

\section{Forces and Freebody Diagrams Extra Reading}
For more information on how to draw free body diagrams and what they can be used for check out this \href{https://www.physicsclassroom.com/class/newtlaws/Lesson-2/Drawing-Free-Body-Diagrams}{website}. The website \href{https://www.phyley.com/}{Phyley} also has more information about free  body diagrams and some worked examples of resolving forces. This includes a version of the mass hanging from two strings problem which sketches how to solve the problem for an arbitrary mass and angle, this examples is found \href{https://www.phyley.com/mass-hanging-from-two-ropes}{here}.

\section{Energy Extra Reading}
The kinetic energy shows up often enough that it may be worth having a formula triangle to help you remember the equation. This is shown in \cref{fig: Kmv triangle}.\\
\begin{figure}[ht]
\ThisAltText{ A formula triangle for kinetic energy, mass, and velocity.}
    \centering
    %\pdftooltip{
    \large \formulatriang[.4]{$E_{K}$}{}{$\frac{m}{2}$}{}{$v^{2}$}{}
    %}{The formula triangle for kinetic energy, mass, and velocity.}
    \caption{The formula triangle for constant velocity, kinetic energy, and mass. }
    \label{fig: Kmv triangle}
\end{figure}

The hyperphysics page about  \href{https://hyperphysics.phy-astr.gsu.edu/hbase/enecon.html}{energy} is a good place to look for extra information.

\section{Circular Motion and Oscillations Extra Reading}
Since the topics in these two weeks are closely related we have grouped them together here. \\

There is a lot more information about circular motion on the Hyper Physics page \href{https://hyperphysics.phy-astr.gsu.edu/hbase/circ.html#circ}{here}. There is also a very interesting YouTube video available \href{https://www.youtube.com/watch?v=AL2Chc6p\_Kk\&ab\_channel=AllThingsPhysics}{here} which talks about what happens to an object undergoing circular motion if the centripetal force is removed. \\

If you want to see a more mathematical description of oscillations, then Chapter 13 of \citep{YandF2019} is where you should look. It explains everything that we discussed in this section of the module, but adds a bit more mathematical sophistication to the discussion. In particular, it shows what happens for the pendulum beyond the simple approximation that we have been working in and explains how to think about energy in a damped system. \\

If you want to learn more about the collapse of the Tacoma bridge then there is a video explaining it \href{https://www.youtube.com/watch?v=3mclp9QmCGs}{here}.

\section{Electric Circuits Extra Reading} 
The webpage \textbf{All About Circuits}, available \href{https://www.allaboutcircuits.com/textbook/}{here}, is full of information about electric circuits that should nicely compliment what we have seen in these lectures. It also has descriptions of lots of experiments that can be carried out to test the basic concepts. Two of these are essentially the same as experiments that you will see in the lab sessions. \\

There are also lots of books about electric circuits in the library. These books are primarily aimed at engineering students and those of you who are going on to do electronic engineering, or who will take Electrical and Electronic Engineering Fundamentals module next year would probably find it useful to familiarise yourselves with at least one of these books.

\section{Modern Physics Extra Reading}
The content of this section is essentially all further reading as it is a non-examinable part of the module which has been included to give you some context about how the topics that you met in this module relate to the current state of physics. Rather than any detailed topics I want to point you towards some of my favourite popular science books. In particular, I want to recommend:

\begin{itemize}
\setlength{\itemsep}{-5pt}
    \item \textbf{Through Two Doors at Once}, \citep{ananthaswamy2020through}, is a wonderful description of the double slit experiment and the profound consequences that it has on how we think about physics at the smallest scales. If you want to know more about quantum physics then this is a great place to start.
    
\item \textbf{Storm in a Teacup}, \citep{czerski2016storm}, this is a great book which shows how to use physics to understand everyday phenomena like how slugs and snails can climb walls and why when you spill coffee it leaves a ring on the table or piece of paper. If you are struggling to know why you should care about physics then this is the book for you. 

\item \textbf{Particle Physics},  \citep{close2023particle}, the whole ``A Very Short Introduction'' series from Oxford University Press is a fantastic place to learn about new topics if you are not sure where to start. The book on particle physics is a great overview of one of the ``big'' disciplines in modern physics where gigantic machines are built to probe the universe at the smallest scale. Particle physics was one the topics that originally got me hooked on physics and I challenge anyone not to find a lot of interesting topics in this book.
\end{itemize}  